% Options for packages loaded elsewhere
\PassOptionsToPackage{unicode}{hyperref}
\PassOptionsToPackage{hyphens}{url}
\documentclass[
]{article}
\usepackage{xcolor}
\usepackage{amsmath,amssymb}
\setcounter{secnumdepth}{-\maxdimen} % remove section numbering
\usepackage{iftex}
\ifPDFTeX
  \usepackage[T1]{fontenc}
  \usepackage[utf8]{inputenc}
  \usepackage{textcomp} % provide euro and other symbols
\else % if luatex or xetex
  \usepackage{unicode-math} % this also loads fontspec
  \defaultfontfeatures{Scale=MatchLowercase}
  \defaultfontfeatures[\rmfamily]{Ligatures=TeX,Scale=1}
\fi
\usepackage{lmodern}
\ifPDFTeX\else
  % xetex/luatex font selection
\fi
% Use upquote if available, for straight quotes in verbatim environments
\IfFileExists{upquote.sty}{\usepackage{upquote}}{}
\IfFileExists{microtype.sty}{% use microtype if available
  \usepackage[]{microtype}
  \UseMicrotypeSet[protrusion]{basicmath} % disable protrusion for tt fonts
}{}
\makeatletter
\@ifundefined{KOMAClassName}{% if non-KOMA class
  \IfFileExists{parskip.sty}{%
    \usepackage{parskip}
  }{% else
    \setlength{\parindent}{0pt}
    \setlength{\parskip}{6pt plus 2pt minus 1pt}}
}{% if KOMA class
  \KOMAoptions{parskip=half}}
\makeatother
\setlength{\emergencystretch}{3em} % prevent overfull lines
\providecommand{\tightlist}{%
  \setlength{\itemsep}{0pt}\setlength{\parskip}{0pt}}
\usepackage{bookmark}
\IfFileExists{xurl.sty}{\usepackage{xurl}}{} % add URL line breaks if available
\urlstyle{same}
\hypersetup{
  hidelinks,
  pdfcreator={LaTeX via pandoc}}

\author{}
\date{}

\begin{document}

{
\setcounter{tocdepth}{3}
\tableofcontents
}
Paul Buy

The Pompeii Project

prehistory

"Decisions were not yet necessary. But the space for them was there."

A dramatic narrative from the Pompeii project

How Michael became who he is -- and why Julia left.

Contents

\hyperref[part-a-boston]{\textbf{PART A -- BOSTON 3}}

\begin{quote}
\hyperref[chapter-a1-last-winter-at-home]{Chapter A1 -- Last Winter at Home 3}

\hyperref[chapter-a2-the-campus]{Chapter A2 -- The Campus 5}

\hyperref[chapter-a3-the-second-language]{Chapter A3 -- The second language 7}

\hyperref[chapter-a4-entry]{Chapter A4 -- Entry 9}
\end{quote}

\hyperref[part-b-rome-preliminary-proximity]{\textbf{PART B -- ROME: PRELIMINARY PROXIMITY 11}}

\begin{quote}
\hyperref[chapter-b1-gregorian-calendar]{Chapter B1 -- The Gregorian 11}

\hyperref[kapitel-b2-living-together]{Chapter B2 -- Living Together 13}

\hyperref[chapter-b3-martina]{Chapter B3 -- Martina 15}

\hyperref[chapter-b4-the-first-trip-in-the-vw-bulli]{Chapter B4 -- The first trip in the VW Bulli 16}
\end{quote}

\hyperref[part-c-the-gentle-break]{\textbf{PART C -- THE GENTLE BREAK 18}}

\begin{quote}
\hyperref[chapter-c1-the-decision]{Chapter C1 -- The Decision 18}

\hyperref[chapter-c2-absence]{Chapter C2 -- Absence 20}
\end{quote}

\hyperref[part-d-pumping]{\textbf{PART D -- PUMPING 21}}

\begin{quote}
\hyperref[chapter-d1-arrival]{Chapter D1 -- Arrival 21}

\hyperref[chapter-d2-childhood-among-ruins]{Chapter D2 -- Childhood Among Ruins 22}

\hyperref[chapter-d3-archaeology]{Chapter D3 -- Archaeology 23}
\end{quote}

\hyperref[part-e-rome-the-silence]{\textbf{PART E -- ROME: THE SILENCE 24}}

\begin{quote}
\hyperref[chapter-e1-the-collegium]{Chapter E1 -- The Collegium 24}

\hyperref[chapter-e2-encounter]{Chapter E2 -- Encounter 25}

\hyperref[chapter-e3-the-silence]{Chapter E3 -- The Silence 26}

\hyperref[chapter-e4-luca]{Chapter E4 -- Luca 27}
\end{quote}

\hyperref[part-f-two-childhoods-parallel]{\textbf{PART F -- TWO CHILDHOODS (Parallel) 28}}

\begin{quote}
\hyperref[chapter-f1-martina-luca]{Chapter F1 -- Martina / Luca 28}

\hyperref[chapter-f2-language]{Chapter F2 -- Language 29}

\hyperref[kapitel-f3-fragment]{Kapitel F3 -- Fragment 30}
\end{quote}

\hyperref[part-g-the-ethician-indirect]{\textbf{PART G -- THE ETHICIAN (INDIRECT) 31}}

\begin{quote}
\hyperref[chapter-g1-the-distance]{Chapter G1 -- The Distance 31}

\hyperref[chapter-g2-a-conversation-on-the-sidelines]{Chapter G2 -- A conversation on the sidelines 32}
\end{quote}

\hyperref[part-h-before-the-workshop]{\textbf{PART H -- BEFORE THE WORKSHOP 33}}

\begin{quote}
\hyperref[chapter-h1-irarah-is-created]{Chapter H1 -- IRARAH is created 33}

\hyperref[chapter-h2-before-the-workshop]{Chapter H2 -- Before the Workshop 34}
\end{quote}

\hyperref[persons]{\textbf{persons 35}}

\begin{quote}
\hyperref[michael-phillips]{Michael Phillips 35}

\hyperref[julia-rossi]{Julia Rossi 35}

\hyperref[martina-rossi]{Martina Rossi 35}

\hyperref[maria]{Maria 36}

\hyperref[luca]{Luca 36}

\hyperref[father-matteo-conti]{Father Matteo Conti 36}
\end{quote}

\section{PART A -- BOSTON}\label{part-a-boston}

\subsection{Chapter A1 -- Last Winter at Home}\label{chapter-a1-last-winter-at-home}

Winter didn\textquotesingle t arrive suddenly; it crept slowly into the days. First the days grew shorter, then the conversations. Finally, snow lay on the streets of Boston, compacted at the intersections, gray at the edges of the sidewalks. The Phillips house stood quietly among them, a narrow row house with a brick facade, nothing special, nothing temporary.

Michael was seventeen, maybe eighteen already, it didn\textquotesingle t really matter anymore. School was behind him like a familiar path, one you could walk even in the dark. In the mornings, he got up early, dressed, ate toast, and drank coffee his mother had already brewed. She was a biology teacher at a high school in the south of the city; her day began in a structured way and rarely ended unexpectedly. His father, a physics teacher, left the house a little later. He read the newspaper thoroughly, as if he needed to understand the world before he left it.

Little was said at the breakfast table. Not out of coldness, but rather out of understanding. They knew each other without constantly saying so. When Michael spoke, both parents listened. When he was silent, no one asked any questions.

Fresh snow had fallen that morning. The road hadn\textquotesingle t been cleared yet, and the cars were parked along the side like half-submerged animals. Michael put on his coat and wrapped his scarf twice around his neck.

"It should stay like this all day," said his mother, placing the lunchbox in front of him.

"Then it will be quiet," said his father, without looking up.

Michael nodded. Silence was not something threatening in this house. It was part of it.

The route to school led past a church, one of many in the neighborhood. He knew the times of the masses, even if he couldn\textquotesingle t recite them from memory. The church was there, like the park or the corner grocery store. You went when it was convenient. You knew it would stay.

In class, Michael usually sat in the second row. Not in the front, not in the back. Physics came easily to him, as did biology. In religion class, he listened, sometimes asking questions, never provocatively. The lessons were about responsibility, about knowledge, about what one should do with what one understood. The answers remained vague, but that didn\textquotesingle t bother him. He had learned that not everything had to be named immediately.

In the afternoon he returned home, took off his shoes, and let the snow melt in the hallway. He sat down at the table in the living room, did his homework, and read. Books were everywhere, not haphazardly, but in small piles. His parents didn\textquotesingle t collect them; they used them.

They ate together in the evening. They had stew, and later tea. The radio was playing softly.

``I spoke to the students about climate models today,'' his mother said. ``Some believe that numbers are opinions.''

His father smiled briefly. "Many people believe that."

"And you?" she asked Michael.

He thought for a moment. "Numbers are tools," he finally said. "But you have to know what they\textquotesingle re for."

His father looked at him. "That\textquotesingle s a good answer."

It contained no praise, more like recognition.

Later, they cleaned up together. Michael dried the dishes, listening to the water in the sink. Through the kitchen window, they could see the street, now empty. The snowflakes fell more heavily, silently.

He sat down at his desk in his room. He had a letter from a university in his drawer, unread. It wasn\textquotesingle t the right moment. Instead, he opened a book and read a few pages, without remembering everything.

He wasn\textquotesingle t thinking about leaving or staying. Only about the fact that something was changing, without him being able to name it. Like the snow that fell and stayed, without being asked.

Late that evening, he heard his parents talking in the living room, their voices hushed and familiar. They spoke about timetables, colleagues, and the coming Sunday. Nothing urgent. That\textquotesingle s precisely why it felt so reliable.

Michael lay down in bed and looked at the ceiling. He felt no restlessness in the sense of resistance. Rather, an anticipation that had not yet taken form. Something that would come because it had to come.

Outside, a snowplow drove past. The sound was short and sharp, then silence again. The road was now clear, the edges higher. Winter had arrived.

Michael closed his eyes. Decisions weren\textquotesingle t necessary yet. But the space for them was there.

\subsection{Chapter A2 -- The Campus}\label{chapter-a2-the-campus}

The campus was elevated, as if one had to leave it slightly to get here. The paths between the buildings were clean, even in autumn when the leaves lay damp on the asphalt. Michael often walked, even if the route was longer. It gave him time to settle in.

He was eighteen, in his first semester. Theology and philosophy were on his timetable, side by side, without explanation. The lecture halls were high-ceilinged, the windows narrow. When the professor spoke, his voice echoed slightly, not disturbingly, more as a way of bringing order to the proceedings. His sentences were clear, his thoughts deliberate. No one had to rush.

Michael usually sat in the middle of the room. He took notes, not everything, just what seemed important to him: names, terms, sometimes a sentence that stuck with him. Other students sat beside him, some quiet, some diligently. During breaks, they stood together, drank coffee from paper cups, and talked about exams, lack of sleep, and the food in the cafeteria.

"You\textquotesingle re from Boston?" someone asked him.

``Yes,'' said Michael.

"Then you\textquotesingle ve already arrived."

Michael smiled. It didn\textquotesingle t feel like it, but he didn\textquotesingle t disagree.

A Jesuit priest led the introductory seminar in the afternoon. He wore a simple black sweater, not a cassock. His voice was calm, almost casual. He asked questions, allowing time for answers. Michael didn\textquotesingle t speak often, but when he did, the others listened.

After the lesson, he sat for a moment. The priest collected the slips of paper and arranged them carefully.

``They are listening attentively,'' he said, without looking up.

Michael wasn\textquotesingle t sure if it was a statement or a compliment. "I\textquotesingle ll try," he said.

The priest nodded. "That\textquotesingle s enough for now."

The library became a fixed place. Afternoons there had their own rhythm. Books were quietly pulled from the shelves, chairs carefully moved. Michael learned to lose himself without giving up. He read texts several times, slowly, sometimes only a few pages at a time.

In the evenings he ate in the cafeteria. The food was simple. Conversations were casual, rarely deep, but honest. Nobody had to prove anything. Studying was work, not a performance.

Sometimes Michael walked across campus after dark. The buildings were lit, some windows still dimly lit. A candle often burned in the chapel. He would sit inside, not regularly, not out of obligation. He simply sat there, listening to the silence.

It wasn\textquotesingle t a place of doubt, but rather one of gathering. Everything was there: knowledge, order, rituals. And yet, something remained open. Michael didn\textquotesingle t feel like a stranger, but neither did he feel like he had arrived. He had the impression that this state was permitted.

In a conversation with a fellow student, he once said: "I don\textquotesingle t really know what I\textquotesingle m looking for here."

The other shrugged. "Same here."

That calmed him more than any answer.

The semesters passed. Exams came and went. Michael passed them without any fuss. He didn\textquotesingle t stand out, but he didn\textquotesingle t disappear either. He was part of the whole, without losing himself in it.

When the first winter arrived on campus, snow covered the lawns. Students walked faster, hunching their shoulders. Michael paused briefly, looking across the square. The campus had become quieter, but not empty.

He knew he was in the right place. Not forever, maybe. But for now.

\subsection{Chapter A3 -- The second language}\label{chapter-a3-the-second-language}

The physics building was closer to the street than the theology faculties. The entrance was austere, glass and concrete, no ornamentation. Michael had to take a different route, a different subway. The ride wasn\textquotesingle t long, but it marked a transition.

He was twenty, maybe twenty-one. His timetable had become more condensed. Lectures in the morning, tutorials in the afternoon, with coffee in between that tasted bitter. The rooms were brighter than at college, more functional. Blackboards covered in equations, which were erased at the end of the day.

In the lab, Michael wore a lab coat that was a little too big for him. He stood at a table with equipment that needed to be precisely calibrated. A lecturer explained the setup, factually and without beating around the bush. Michael listened, asking questions when something remained unclear. No one seemed surprised that he was also studying theology. It was noted, but not commented on.

"You\textquotesingle re from college, right?" asked a fellow student as they connected cables.

``Yes.''

"Then you\textquotesingle re used to thinking that way."

Michael smiled. "Maybe a different one."

They continued working in silence. Numbers appeared on the screen, were noted down, checked. Mistakes were part of the process, nothing personal. If something didn\textquotesingle t work, they started over.

In the lectures, Michael sat further forward. He wanted to see how the formulas were developed. The language was concise and precise. There was little room for embellishment. A single sentence could contain everything if it was correctly formulated.

After class, he sometimes lingered in the hallway with other students. Conversations revolved around assignments, results, and the next internship. The pace was different, faster, more intense. Michael adapted without changing himself.

That evening, back in his apartment, he laid the notebooks side by side. Notes from theology, from physics. He didn\textquotesingle t mix them. It was important to him that each language had its own space. He felt no contradiction, rather an expansion.

One day after a seminar, a lecturer stopped briefly to talk to him. "You work meticulously," he said. "You take your time."

Michael nodded. "I need to understand what I\textquotesingle m doing."

The lecturer looked at him, scrutinizing him, then approvingly. "That\textquotesingle s rare."

Michael often walked home. The city was loud, even in the evenings. Sirens, voices, traffic. He let the sounds wash over him. Formulas, concepts, observations began to take shape in his mind. It wasn\textquotesingle t a compulsion, more a gradual merging.

Sometimes he would bump into fellow college students on the street. They would greet each other, exchange a few words. Their worlds touched without merging.

Michael felt no inner conflict. He learned that you didn\textquotesingle t have to say everything at once. Some things required precision, others patience. Both could coexist.

As the semester ended, he stood in the lab one last time, looking at the equipment, now switched off. He had learned something that went beyond the formulas. A second language he understood.

He left the building and stepped out into the evening. The air was cool and clear. Boston lay before him, familiar yet open.

\subsection{Chapter A4 -- Entry}\label{chapter-a4-entry}

The decision didn\textquotesingle t come on a specific day. There was no moment when one could say: Now. It had been lying dormant in Michael\textquotesingle s everyday life for some time, like an object that one puts aside and touches again and again to make sure it\textquotesingle s still there.

Michael was twenty-two, perhaps already twenty-three. He had organized his studies, planned his degrees, and had his options open. In his conversations with the Jesuits, there was no pressure. They asked questions, listened, and let time pass.

A priest received him in a small office on the edge of campus. There were no stacks of files on the desk, only a notebook and a pen.

"They stayed for a long time," said the priest.

Michael nodded. "I wanted to be sure."

"Certain about what?"

Michael thought for a moment. "That I don\textquotesingle t flee."

The priest smiled slightly. "That\textquotesingle s a good start."

More conversations followed. They talked about daily life, discipline, and community living. The discussion focused on rules, not visions. Michael listened attentively and asked questions. He wanted to know what days were like, how years passed.

He told his parents about it one evening after dinner. The table had been cleared, but the tea was still hot.

``I will join the order,'' he said.

His mother looked at him calmly. "When?"

"In autumn."

His father nodded slowly. "Then you took your time."

``Yes.''

Nothing more was said. There were no warnings, no explanations. They knew that Michael didn\textquotesingle t make decisions to prove anything.

The preparations were unspectacular. Forms, conversations, appointments. Michael organized his books, gave some away, kept others. He said goodbye to fellow students without much fanfare. Some stopped and asked questions. Others simply took it in, as if it were a change of location.

The day of the initiation was cool. A small room, few people. A simple ritual, clear words. Michael spoke them calmly, without hesitation. It didn\textquotesingle t feel like an ending, more like a transition.

After the service, they drank coffee. They talked about organizational matters, about the next step. Rome was mentioned, casually, as something that was to come.

That evening, Michael walked through the city one last time. The streets were familiar, the sounds the same. He knew he would soon be leaving. It didn\textquotesingle t make him sad.

He stood just outside his parents\textquotesingle{} house and looked up at the window. The light was on. He went inside and sat down next to it. Nothing had been left open.

The order he had chosen was not an escape. It was a framework. And within that framework, he remained the same.

\section{PART B -- ROME: PRELIMINARY PROXIMITY}\label{part-b-rome-preliminary-proximity}

\subsection{\texorpdfstring{Chapter B1 --\emph{Gregorian calendar}}{Chapter B1 --Gregorian calendar}}\label{chapter-b1-gregorian-calendar}

The Gregorian University was bathed in partial shade that morning. The sun only reached the courtyard at an angle, and the students\textquotesingle{} voices quickly faded among the walls. Michael was new to the city, still unsure of his way around. He usually arrived a few minutes early, sat in the second row, and neatly arranged his notes in front of him.

Julia sat two rows ahead. She had dark hair, which she usually wore tied back, not intentionally, but rather out of habit. Michael barely noticed her at first. Only when the lecturer asked a question and she answered did he look up. Her voice was calm, clear, and unhurried. She spoke Italian, then switched effortlessly to English.

After the seminar, some students remained standing in the room. Michael gathered his things, and Julia stood beside him.

"You\textquotesingle re not from here," she said in English.

``Boston,'' he replied.

"I thought to myself, this isn\textquotesingle t Europe."

They both smiled briefly.

Over the next few weeks, they saw each other regularly. In the library, at the tables in the back, where it was quieter. They worked side by side, occasionally exchanging books. Conversations arose casually, about texts, about concepts, about the pace of the lectures.

Once they were sitting in a small café opposite the university. The coffee was strong, the table wobbled.

"The seminars here are slower than in the USA," said Michael.

``We talk more about it,'' Julia replied.

"Maybe," he said. "Sometimes that helps."

She laughed softly. "Not always."

They talked about their journeys to Rome. Julia had come from northern Italy, having previously studied in Milan. Michael spoke of college, of switching between disciplines. It didn\textquotesingle t sound like a justification, more like a description.

They didn\textquotesingle t meet intentionally, but they did so again and again. Sometimes they walked together for a while after leaving the library. They talked about little things: the weather, the noise of the city, the cramped conditions of their apartments. Their closeness didn\textquotesingle t arise from gestures, but from reliability.

One evening they sat on the steps near the Piazza della Pilotta. The city was loud, the traffic heavy. Michael watched the people passing by.

"Rome is exhausting," he said.

Julia nodded. "But something remains."

He knew what she meant, even though she didn\textquotesingle t explain it.

There was no single decision, no standout moment. They began to coordinate their work, sending each other texts, passing on questions. If one was missing, the other noticed.

The relationship didn\textquotesingle t begin as such. It was first a collaboration. Then a habit. Only later did it get a name.

For the moment, it was enough to sit and work next to each other.

\subsection{\texorpdfstring{Kapitel B2 --\emph{living together}}{Kapitel B2 --living together}}\label{kapitel-b2-living-together}

The apartment was on the third floor of a building near Via Merulana. Julia had lived there since her first semester. Two rooms, a narrow kitchen, windows facing the courtyard. The plaster was crumbling in places, and in winter there were drafts through the frames. She had gotten used to it.

Michael came almost every evening. He knew the way, the steps, the moment when the light in the stairwell flickered. He usually carried a bag with books, sometimes bread or fruit from the shop on the corner. He didn\textquotesingle t ring the bell for long. Julia opened the door, let him in, and asked no questions.

They cooked together. Simple food, nothing elaborate. While the water heated up, one of them talked about their day. About a seminar, a casual remark, a statistic that didn\textquotesingle t add up. Michael spoke calmly, Julia listened, added to it, occasionally disagreed.

"You haven\textquotesingle t eaten enough again today," she said once, and placed a plate in front of him.

``I forgot,'' he replied.

They sat down at the table, side by side. The books mostly remained open, with notes scattered among them. Sometimes one of them read a paragraph aloud. If something remained unclear, they left it as it was.

Michael went back later. He didn\textquotesingle t always say goodbye with words. Sometimes he just grabbed his bag and put on his coat. Julia would stand by the door, listening to his footsteps disappear into the stairwell.

The Jesuit house was quiet. The corridors were long, the doors closed. Michael placed the books on the desk and jotted down a few more thoughts. The day wasn\textquotesingle t over, but it was organized.

The next morning they saw each other again at the Gregorian University. Briefly, casually. A nod, a smile. They didn\textquotesingle t sit next to each other, but they knew where the other was.

They stayed together longer over the weekend. They walked through the city, seeking leisurely routes. Small squares, narrow streets. They talked about Rome, about staying and moving on, without making any firm decisions.

``We are here for a while,'' Julia once said.

``Yes,'' said Michael.

It wasn\textquotesingle t a farewell, just an observation.

Church was part of the routine. Michael went at the set times; Julia sometimes went with him, sometimes not. There were no explanations for this. It wasn\textquotesingle t an issue that needed to be decided.

In the evenings they sat in the kitchen, the window open, voices drifting from the courtyard. They said little. Closeness wasn\textquotesingle t shown through gestures, but in the fact that neither left as long as the other was still there.

It wasn\textquotesingle t a hidden life. Nor was it a public one. It was simply their everyday life, between university, city, and the paths they walked anew each day.

\subsection{Chapter B3 -- Martina}\label{chapter-b3-martina}

The hospital was located just outside the city center. The route there led through streets Michael knew, but they felt different this morning. The traffic was heavy, the light gray. Julia sat beside him, silent, her hands resting on her stomach. She said little. When she did speak, her speech was clear.

The delivery room was brightly lit. Voices, footsteps, brief instructions. Michael stood to the side, holding Julia\textquotesingle s hand, letting go when necessary. He did as he was told, without question. Julia breathed calmly, with concentration. She didn\textquotesingle t seem anxious, but rather composed.

When Martina was born, there was a brief silence. Then a sound, surprisingly strong. Julia closed her eyes and exhaled. Michael saw the child, small, curled up, her skin still red. He didn\textquotesingle t know what to feel. He only knew that she was there.

The midwife laid Martina on Julia\textquotesingle s chest. Julia put her arm around her as if she\textquotesingle d always done so. Michael stepped closer and looked down at them both. Nobody said anything.

The first few days were chaotic. Time blurred. Michael came in the morning, left in the evening, and came back again. He learned to recognize the movements, the brief crying that subsided again. Julia rested, got up, and sat down again. She accepted everything without complaint.

Little changed in the apartment, and yet everything did. A small bed in the bedroom, blankets, bottles. The table remained the same, but people sat at it differently. Conversations became shorter, glances longer.

Michael held Martina alone for the first time, uncertainly. Julia watched him, without correcting him.

"That\textquotesingle s good," she said.

He nodded, stopped, and barely moved.

A rhythm gradually settled in. Sleep came in segments. Food became secondary. Michael adjusted his schedule. He arrived earlier, sometimes stayed longer. He left again when necessary.

Outside, life went on. Lectures, appointments, street noise. Inside, something had emerged that needed no explanation.

There were no conversations about the future or order. Only the present. Martina slept, woke up, slept again. Julia watched her attentively. Michael did too.

In the evening they sat together, the light dimmed. Martina lay between them, breathing calmly. No one spoke. Nothing had been decided. But something was there, something that remained.

\subsection{Chapter B4 -- The first trip in the VW Bulli}\label{chapter-b4-the-first-trip-in-the-vw-bulli}

The VW bus was old but reliable. Light blue, with small rust spots along the edges. Michael had taken it over from a fellow student who was leaving Rome. The engine sounded rough but steady. Julia liked that sound. It promised movement, without haste.

They set off early. The city was still quiet, the streets emptier than usual. Martina was asleep in the back seat, wrapped in a blanket. Her breathing was calm. Julia kept turning to check on her, making sure everything was alright.

"She is sleeping soundly," she said.

Michael nodded, keeping his eyes on the road. "She likes that."

They drove towards the sea, without a fixed destination. The windows were open, warm air came in. Michael turned on the radio and let it play softly. Music that neither of them recognized.

They stopped at a small square. Trees provided shade. Michael parked the car and turned off the engine. The silence that followed was almost palpable.

They spread out a blanket. Julia sat down and picked Martina up. The child briefly opened her eyes, then closed them again. Michael sat down beside her and handed Julia a bottle of water.

"This is good," said Julia.

``Yes,'' said Michael.

They ate bread, cheese, and fruit. Nothing special. The van was parked behind them, like a shield. Every now and then someone walked by, nodded, and said nothing.

In the afternoon they went for a walk. Michael carried Martina, carefully, inexperienced. Julia walked beside him, briefly holding his shoulder when the path became uneven.

"Slowly," she said softly.

He smiled, adjusted his pace.

Later, when it got cooler, they continued driving. Martina fell asleep again. The sun was low in the sky. Michael drove steadily, without accelerating. Julia leaned her head against the window and looked out.

There were no discussions about what was to come. No plan, no promises. Just the path that opened up before them.

That evening they parked at a basic campsite. Michael parked the van and opened the door. The air smelled of grass and dust. Julia laid Martina in the van and adjusted the blanket.

They sat outside for a while longer, side by side. The light grew dim. Martina slept peacefully.

It was nothing extraordinary. But it was theirs.

\section{PART C -- THE GENTLE BREAK}\label{part-c-the-gentle-break}

\subsection{\texorpdfstring{Chapter C1 --\emph{The decision}}{Chapter C1 --The decision}}\label{chapter-c1-the-decision}

They were sitting in the kitchen. It was late evening, the window was open. Sounds from the courtyard drifted in: voices, footsteps, the clatter of dishes. Martina was asleep in the next room.

Julia had a letter lying in front of her, unopened. She ran her finger over the envelope without lifting it. Michael sat opposite her, his hands around a cup that was long since cold.

``I spoke to someone today,'' said Julia.

Michael nodded and waited.

"In Pompeii. A position. At the family counseling center."

She looked at him, calmly, without expectation.

"When?" he asked.

"In autumn."

There was a brief silence. Michael looked towards the window, then back at her. "That\textquotesingle s far," he finally said.

``Yes.''

She opened the envelope, pulled out the paper, and put it back down. "It\textquotesingle s good work. And I could get closer to..." She broke off, looking towards the next room.

Michael understood. He said nothing.

``I can't stay,'' she continued. ``Not like this.''

"I know," he said quietly.

He had known it long before she said it. Not as a thought, but rather as a certainty that had slowly settled in.

``I don't want it to be\ldots'' she began, then paused. ``To be harder than it needs to be.''

Michael nodded. "Me neither."

They looked at each other. There was no bitterness between them, more like weariness. The kind of weariness that comes from carrying something for a long time.

``We will stay connected,'' said Julia.

``Yes,'' said Michael.

"For Martina."

``For her,'' he confirmed.

He stood up, walked a few steps across the room, and stopped. "I\textquotesingle ll come," he said. "As often as I can."

Julia smiled weakly. "I know."

They stood there in silence for a while. Then Julia went into the bedroom and looked for Martina. Michael didn\textquotesingle t follow her. He stayed in the kitchen and put away the cups.

It wasn\textquotesingle t a separation with doors slamming shut. It was a step that had to be taken.

Later, as he left, he took his jacket from the chair. Julia was standing by the door. They hugged each other briefly, firmly, without haste.

"Take care of yourself," she said.

"You too."

The door closed quietly.

\subsection{Chapter C2 -- Absence}\label{chapter-c2-absence}

After Julia left, Michael\textquotesingle s daily routine became clearer. He got up early, walked his usual routes, and returned home in the evening. The apartment was no longer part of his everyday life. Instead, he spent more time in the library, in the Gregorian University\textquotesingle s study rooms, and in quiet offices with tall shelves.

He set up a permanent workspace. The desk stood near the window, allowing the light to stream in evenly. Books lay open, notes neatly stacked. He worked with concentration, rarely interrupting himself. When he left the room, he left everything as if he were about to return.

Conversations with colleagues became shorter. They talked about texts, appointments, deadlines. Michael was reliable. He came prepared, listened, and asked questions when necessary. No one asked about his private life. There was no reason to.

In the evenings, he ate in the refectory. Seats changed, as did faces. Conversations were calm, often matter-of-fact. Michael participated without being intrusive. When it became quiet, he remained seated until it was time to leave.

Sometimes he wrote letters. Short sentences, factual. He asked about Martina, about Julia\textquotesingle s work. The replies came irregularly, but they came. He put the letters in a drawer, neatly folded.

The workload increased. The inventory project demanded attention. Feedback had to be evaluated, adjustments made. Michael worked carefully, step by step. It wasn\textquotesingle t an escape. It was what was possible at the time.

He often sat at his desk late into the night. The corridors were silent, with only occasional lights flickering behind doors. He wrote, erased, and rewritten. Not out of restlessness, but out of meticulousness.

Sometimes he went for walks on weekends. Without a destination. The city had become familiar, without feeling close to him. He walked past cafes, past places where he used to sit. He didn\textquotesingle t stop.

There was no bitterness. Only distance. Things shifted, slowly, silently.

Michael noticed something was loosening. Not suddenly, but like a knot untying itself over time. The connection remained, but it no longer provided support.

He continued working. And that was enough.

\section{PART D -- PUMPING}\label{part-d-pumping}

\subsection{\texorpdfstring{Chapter D1 --\emph{Arrival}}{Chapter D1 --Arrival}}\label{chapter-d1-arrival}

The train stopped abruptly. Julia took the bag from the luggage rack, set it down briefly, and picked Martina up. It was warm outside, despite the early hour. The air smelled different than in Rome---thickener, sweeter, mixed with dust.

The platform was narrow. Voices echoed between the walls. A man shouted something Julia didn\textquotesingle t understand, but it didn\textquotesingle t sound important. She walked slowly, watching her steps. Martina held onto her shoulder and looked around, still, alert.

The apartment was on a side street, not far from the center. Two rooms, a small kitchen. The shutters were closed. Julia opened them, one after the other. Light streamed in, bright, unrestricted. You could hear cars, voices, music somewhere.

She put the suitcase down and sat Martina on the floor. The child took a few steps, stopped, and sat down. Julia laughed softly. She took a cloth from her bag, wiped the table, and opened a cupboard. It smelled of dampness and cleaning products.

In the afternoon, she went out. The walk to the supermarket was short. She bought bread, water, and some fruit. Martina sat in the buggy, her eyes half-closed. The street was noisy, but not hectic. People stopped, talked, and then continued on their way.

In the evening they ate at the small table. Julia cut the bread and handed Martina a piece. Outside it grew quieter. A dog barked, then nothing more.

Later, she put Martina to bed. The room was still warm. Julia sat down on the chair next to her and stayed there for a while. It wasn\textquotesingle t a goodbye or a beginning. Just a place that was now hers.

She would go to work the next morning. Not today. Today it was enough to have arrived.

\subsection{Chapter D2 -- Childhood Among Ruins}\label{chapter-d2-childhood-among-ruins}

Martina knew the paths among the ruins before she could read. Julia often took her along when she made a detour after work. They walked slowly, stopping frequently. Martina would sit on low walls and run her fingers over the stone. It was warm.

The school was on the outskirts of town. Julia took Martina there in the mornings and picked her up in the afternoons. Martina carried a backpack that was almost too big. The classroom smelled of chalk and dust. The teacher spoke loudly, the children all at once. Martina listened, rarely raising her hand.

In the afternoon they played outside. Tourists walked by, holding up cameras, speaking foreign languages. Martina paid them no attention. For her, the ruins weren\textquotesingle t a destination, but rather her surroundings: walls, shadows, stairs that led nowhere.

At home, there was a small table where Julia worked. Files, notebooks, sometimes a book. Martina sat beside her, drawing or doing homework. When she asked where her father was, Julia answered calmly. He lived in Rome and worked a lot. She said nothing more. That was enough.

Sometimes mail arrived from Rome. Envelopes, neatly addressed. Julia put them aside, opening them later. Martina didn\textquotesingle t ask about them. The father wasn\textquotesingle t a secret, but neither was he a part of everyday life.

On their days off, they set off early. Through the streets, past the closed shops, into the compound. The city was quiet then. Martina walked ahead, turned around, and waited. Julia followed.

In the evening they sat on the balcony. Sounds from below, voices, dishes. Martina talked about school, about a wall she had discovered that day. Julia listened. The sky grew dark, slowly.

It wasn\textquotesingle t a special life. But it was hers.

\subsection{Chapter D3 -- Archaeology}\label{chapter-d3-archaeology}

Martina lived in a small room near San Lorenzo. The walk to the university was short, past bars that opened early. In the mornings, the air smelled of coffee and exhaust fumes. She usually walked alone, with a backpack containing notebooks, books, and an old measuring tape.

The Sapienza buildings seemed austere. Long corridors, stone staircases, lecture halls with worn seats. Martina often sat in the middle rows. She took notes, slowly, neatly. Names, dates, shift sequences. Nothing about it was spectacular. It was work.

She spoke little in the seminars. She listened, asking questions when she was sure of her abilities. Her lecturers soon got to know her---not as particularly loud, but as reliable. If something needed to be submitted, she did. If an excavation was planned, she volunteered.

Her first excavation took her to the area surrounding Rome. Early mornings, long days. Dirt under her fingernails, dust in her hair. She learned to distinguish layers, to work carefully, not to rush anything. In the evenings, they sat together, ate simple food, and talked about their finds. Martina listened, asked questions, and took notes.

Sometimes she called her mother. Short conversations. Julia asked about her studies, about her daily life. Martina talked about the heat, her bones, the walls that were slowly becoming visible. They didn\textquotesingle t talk much about personal matters. It wasn\textquotesingle t necessary.

In her second year, she was given more responsibility. She was allowed to manage a small area on her own. She marked, documented, and photographed. In the evenings, she organized the data, compared plans, and made corrections. She liked that nothing happened quickly.

In between, she traveled to Pompeii. She knew the streets, the ruins, the rhythm of the city. She walked through it with new eyes, pausing where she had previously simply continued on. History was now something concrete.

Back in Rome, in her room, she often sat at her desk with the window open. Sounds from outside: voices, scooters. She read, wrote, crossed out sentences. There were days when she was tired, and days when she was sure she was where she belonged.

At the end of her studies, she submitted her thesis. It wasn\textquotesingle t a big moment. She handed in her folder, went outside, and sat down on a bench. The sun was high in the sky. She didn\textquotesingle t think about what was to come. It was enough that she had found her place.

\section{PART E -- ROME: THE SILENCE}\label{part-e-rome-the-silence}

\subsection{\texorpdfstring{Chapter E1 --\emph{The Collegium}}{Chapter E1 --The Collegium}}\label{chapter-e1-the-collegium}

\emph{Maria arrived early. She unlocked the office door, put down her bag, and raised the blinds. The desk was tidy, as always. Calendars, filing trays, a stack of letters.}

\emph{She made coffee in the small kitchen at the end of the hall. While she did, she heard footsteps on the stairs, muffled voices. The seminarians were going to breakfast. Maria knew many of them by name, knew who was on time and who wasn\textquotesingle t. She nodded to them, nothing more.}

\emph{In the office, she sorted mail, typed up appointments, and prepared files. German, Hungarian, Italian. Sometimes English. She switched effortlessly. It was part of the job. The institution was old, the procedures established. Maria fit in without attracting attention.}

\emph{Around noon, she carried documents to another floor. She walked slowly, careful not to disturb anyone. Doors were open, and soft conversation drifted from a room. She didn\textquotesingle t stop.}

\emph{During the break, she sat briefly in the courtyard. The sun had finally reached the stones. She ate a sandwich and watched the birds. The place was sheltered, enclosed. Outside, the city was noisy, but not here.}

\emph{Visitors arrived in the afternoon: priests, professors, and guests from abroad. Maria greeted them, took their coats, and showed them the way. She was friendly and reserved. She asked no questions.}

\emph{In the evening, she locked the office. She walked through the rooms one last time, checking the windows and doors. Everything was in its place. On her way out, she paused briefly and looked back. The Collegium lay still, unchanged.}

\emph{It wasn\textquotesingle t her place, but she worked there. That was enough.}

\subsection{Chapter E2 -- Encounter}\label{chapter-e2-encounter}

They first saw each other in the corridors. Michael was coming from the reading room, Maria was carrying a stack of files. They nodded to each other, as was customary here. Later, they exchanged a few words. Appointments, a form, a signature.

Sometimes they met in the courtyard. Maria sat on the bench, Michael walked by, paused briefly. They talked about the weather, about the noise of the city outside. Nothing personal. That was enough.

Once, Maria brought him some documents to his room. Books lay on the table, papers beside them. She set the envelope down and waited. He thanked her and asked if she had been working there long. She answered briefly. He nodded.

They often met in the kitchenette. Michael stood by the window, Maria filled cups. They exchanged a few words. He asked for a phone number, she gave it to him. It was for a follow-up question, she said. He understood.

The connection was there, but it remained quiet. No unnecessary gestures, no superfluous words. They both knew where they were. The institution set boundaries without explicitly stating them.

When they said goodbye, it happened without any particular reason. Michael walked faster, Maria paused for a moment. Then she too continued on her way.

\subsection{\texorpdfstring{Chapter E3 --\emph{The silence}}{Chapter E3 --The silence}}\label{chapter-e3-the-silence}

Maria noticed the change early on. Not as a shock, more as a shift. She became more tired, paid closer attention to her body. At the office, she continued working as usual, organizing files, answering calls. Nobody questioned her.

She went to her appointments alone. The waiting rooms were quiet and businesslike. She listened, nodded, and signed. On her way home, she bought bread, fruit, and water. Her daily routine continued.

Michael continued to see her. In the hallways, in the courtyard. He noticed that she stopped less often, that she walked faster. He didn\textquotesingle t ask any questions. Maria was grateful for that.

When her belly became visible, she changed her clothes. More jackets, more routes. She took time off work earlier, said she was fine. And she was.

The birth was brief. A room, white light, a voice saying her name. Then the child lay beside her, small, warm. She held him tight, as if she needed to protect him.

Luca slept a lot. Maria watched him, learned his movements. She gave him a name without thinking about it for long. It wasn\textquotesingle t a promise, just a beginning.

She thought of Michael, sometimes. Of his calm demeanor, of his hands when he spoke. She also thought of the house, the rules, the paths laid out for him. She imagined what one word would change. Too much.

So she remained silent. Not out of fear, not out of defiance. She wanted to protect the child, and him too. The silence was a choice, not a deficiency.

As she walked home with Luca, she carried him close to her body. The city was loud, the sky clear. Maria walked slowly. She had time.

\subsection{Chapter E4 -- Luca}\label{chapter-e4-luca}

Luca grew up in a small apartment not far from the city center. The windows faced a narrow street, with parked scooters below. Maria placed his bed by the window to let in fresh air. He slept peacefully.

She took him to daycare early. In the mornings they went together, and in the afternoons she picked him up. Luca didn\textquotesingle t talk much, but observed a lot. He liked to sit on the sidelines, build with blocks, and sort them by color. Maria let him.

There were books at home. Old editions, textbooks, calendars. Luca would leaf through them before he could read. Maria would explain what he asked. Nothing more.

He didn\textquotesingle t stand out at school. He was neither particularly good nor particularly bad. The teachers said he was quiet and reliable. Maria nodded. She knew him differently.

Luca often stayed alone in the afternoons. Maria worked longer hours and didn\textquotesingle t return until evening. Luca would make himself something to eat and turn on the computer. He learned how things worked, not out of ambition, but because it calmed him.

They sometimes went to the Collegium. Maria still had contacts there and would bring documents by. Luca waited in the courtyard, looking at the fountain and the walls. He knew the place, even though he didn\textquotesingle t belong there.

He had few friends. One boy from the neighborhood, and later no one else. Luca wasn\textquotesingle t unhappy, but he kept to himself. He asked questions that Maria couldn\textquotesingle t always answer. Sometimes she said nothing. That, too, was an answer.

At sixteen, he had grown tall, thin, and attentive. He could take care of himself and find his way around. Maria knew he would go his own way. She accompanied him as long as she could.

\section{PART F -- TWO CHILDHOODS (Parallel)}\label{part-f-two-childhoods-parallel}

\subsection{\texorpdfstring{Chapter F1 --\emph{Martina / Luca}}{Chapter F1 --Martina / Luca}}\label{chapter-f1-martina-luca}

Martina sometimes came back to Pompeii on weekends. The apartment was still the same. Julia cooked, set the table, and asked about her studies. Martina told her about seminars, an excavation outside Rome, and long days. It was a calm, unhurried way of talking.

During the week, she lived in Rome. Sapienza University had its own rhythm. In the mornings, the streets filled up; by midday, they emptied again. Martina sat in the library, working her way through texts, making sketches. She knew the faces around her. People greeted each other, sat down next to each other, and continued working.

In the evenings, she walked home. The streets were loud, familiar. She bought bread, sometimes some cheese. At home, she continued writing, preparing for the next day. Her life was full, but manageable.

Luca moved around the same city, but differently. He didn\textquotesingle t know the fixed schedules. He knew when libraries closed and which ones stayed open. He sat at tables where no one asked questions. He used computers that weren\textquotesingle t registered to him.

During the day, he was inconspicuous. He walked through neighborhoods where you didn\textquotesingle t stand out. He overheard conversations, memorized names. In the evenings, the city became quieter. Luca worked then. He wrote, deleted, started again.

Sometimes he walked past the Collegium. He stayed outside and sat on a wall. Inside, he knew the paths, but they didn\textquotesingle t belong to him. Maria still worked there. They rarely asked each other questions.

Martina took the metro, standing crowded among other students. Luca avoided her. He walked, taking detours. Both knew the city, both moved around confidently. But they used different doors.

One evening, Luca was sitting in an internet café. Someone next to him was talking about Milan, about contacts, about a trip. Luca listened. He said nothing. Rome was still there, but not enough anymore.

Martina was sitting at her desk at the same time. She was organizing notes, putting books aside. Cars drove by outside. She was thinking about the next day.

They didn\textquotesingle t know each other. But the city held them both, for a time.

\subsection{Chapter F2 -- Language}\label{chapter-f2-language}

Martina learned to speak precisely from an early age. At the school in Pompeii, this was emphasized. The teachers expected complete sentences, neat notebooks, and punctual submissions. Julia sat at the table in the evenings, helping with homework, proofreading texts, underlining nothing, and only asking questions. Martina answered, sometimes briefly, sometimes at length. Both were acceptable.

Italian was spoken at home, English was added at school, and later Latin. Words had their place. They weren\textquotesingle t used carelessly. Martina memorized phrases and listened for transitions. She liked it when things sounded right.

When she went to Rome, the atmosphere changed. At the university, language moved faster. Lecturers spoke ahead, tacitly assuming a certain level of knowledge. Martina sat in the lecture halls, taking notes and marking passages she would later have to review. Seminars were filled with discussion. She didn\textquotesingle t often participate, but when she did, she was well-prepared.

A different order prevailed in the library. Books were put back, workspaces kept clear. Martina soon knew the floors, the quiet tables, the best times to arrive or leave. She laid her materials side by side and put them back in their proper place at the end of the day.

During excavations, she learned a different language. Short instructions, clear agreements. They pointed out find spots, measured, and took notes. Dust hung in the air, voices were hushed. Martina worked with concentration, listened, and did what was necessary.

In the evenings, she wrote reports. She paid attention to wording and adhered to guidelines. The university demanded clarity, not embellishment. Martina was comfortable with this. She knew what was expected.

Institutions were not foreign to her. School, university, excavation teams -- they provided structure. Martina moved confidently within them. She knew the rules and used them without giving them much thought.

Language was part of it. It carried her through the day, through texts, conversations, and minutes. It organized, connected, and preserved. Martina didn\textquotesingle t talk much about herself. But what she did say endured.

\subsection{Kapitel F3 -- Fragment}\label{kapitel-f3-fragment}

Luca grew up in apartments that never seemed quite finished. Sometimes a shelf was missing, sometimes a curtain. Maria made sure it was clean. More wasn\textquotesingle t always possible. She worked long hours, came home tired in the evenings, dropped off groceries, and asked about school. Luca answered briefly. He had learned to entertain himself.

He changed schools frequently. Not often, but often enough that friendships remained tenuous. Luca usually sat in the back, observing. He wasn\textquotesingle t disruptive. He stood out more because he knew things that weren\textquotesingle t part of the curriculum: number sequences, terms, program names. The teachers took note of it without asking many questions.

He spent his afternoons outdoors or in front of screens. Internet cafes, later secondhand laptops. He learned where to find power outlets, how to buy time without spending much money. He read forums, manuals, and fragments of texts. He jumped between topics, rarely staying with one for long. He saved what interested him.

Maria didn\textquotesingle t ask many questions. She made sure he got up in the morning and ate in the evening. Sometimes they sat together and watched the news. Luca listened, offering brief comments. For him, politics was something that ran in the background, like traffic noise.

He rarely went to church. He only knew the college from the outside. Large doors, clear procedures. Maria didn\textquotesingle t talk about it. Luca didn\textquotesingle t ask any questions.

His education was piecemeal. An online course here, an abandoned project there. He learned quickly, but forgot some things again. He knew how to gain access, how to organize information. Certificates didn\textquotesingle t play a major role.

He was often awake at night. He wrote snippets of code, read chats, followed discussions that had already disappeared the next day. Contacts were made and dissolved. Names remained vague.

Luca wasn\textquotesingle t lost, but he wasn\textquotesingle t integrated either. Rome offered him spaces where he could stay without attracting attention. He moved confidently within them. His world consisted of transitions, open ends, fragments.

\section{PART G -- THE ETHICIAN (INDIRECT)}\label{part-g-the-ethician-indirect}

\subsection{\texorpdfstring{Chapter G1 --\emph{The distance}}{Chapter G1 --The distance}}\label{chapter-g1-the-distance}

\emph{The news didn\textquotesingle t come officially. No circular, no meeting. Michael heard about it between appointments, in the hallway of the Gregorian University, where voices remained hushed, even when no one was listening.}

\emph{"Conti\textquotesingle s gone," a colleague said, almost casually. He was holding a folder under his arm, as if he were in a hurry. Michael stopped. "Gone?" -- "Transferred. Retired. Something like that."}

\emph{In the following days, Conti was simply missed. His seat in the seminar room remained empty. The door to his office was closed, the nameplate still in place. Students waited in vain after the lecture, as they usually did. No one explained anything.}

\emph{Michael remembered conversations that never lasted long. Conti spoke softly, asked questions that remained open. He never gave instructions, never defended positions. Sometimes he simply listened, nodded, jotted down a sentence. Michael had taken this for granted.}

\emph{Now the tone changed. Other lecturers took over topics, speaking in a more structured, clearer way. There were more references, more reassurances. The seminars continued to work. Everything worked.}

\emph{Once, Michael Conti met by chance near the Piazza della Pilotta. He wasn\textquotesingle t wearing a gown, just a coat. They paused briefly. Conti asked about his work, about the students. Michael answered matter-of-factly. There was no parting remark.}

\emph{As they parted, Conti simply said: "Stay alert." Not emphatically, not admonishingly. Like a casual remark.}

\emph{Michael went back to the university. The day continued as if nothing had happened. But something had shifted. Invisibly, indefinably. More like a subtle loss of weight, only noticeable when lifting something that used to be heavier.}

\subsection{Chapter G2 -- A conversation on the sidelines}\label{chapter-g2-a-conversation-on-the-sidelines}

The room was at the back of the building. There was no sign on the door, just a narrow corridor that ended there. Luca had come for the files. Someone had told him they were there, disorganized, hardly used. He had nodded, asked nothing further.

Conti sat at a table near the window. The light was subdued, the afternoon well advanced. In front of him lay a stack of papers, handwritten notes, nothing official. He looked up as Luca entered, not surprised, more attentively.

"Are you looking for something?" he asked.

Luca mentioned the name, the time period. Conti listened, slowly stood up, and went to a shelf. He pulled out a folder and placed it on the table. His movements were calm, without haste.

Luca sat down and turned the pages. Conti remained standing and looked out the window. Muffled noise drifted in from outside---footsteps, voices, nothing in particular.

"Do you work here?" Conti asked after a while.

``Not exactly,'' said Luca. He said nothing more.

Conti nodded. "Direct communication is rarely necessary," he said, almost casually.

They were silent. Luca jotted something down and pushed the folder back. Conti sat down again and clasped his hands. He asked another question, about Luca\textquotesingle s studies, about what interested him. Luca answered briefly and matter-of-factly. It didn\textquotesingle t sound evasive, more cautious.

Conti listened without interrupting. Finally, he said: "Pay attention to what you can\textquotesingle t categorize. That often lingers longer."

Luca glanced up briefly. He said nothing. He stood up and thanked him. Conti returned the greeting with a slight nod.

When Luca left the room, Conti remained seated. He picked up his notes again, finishing a sentence as if the interruption had never happened.

\section{PART H -- BEFORE THE WORKSHOP}\label{part-h-before-the-workshop}

\subsection{\texorpdfstring{Chapter H1 --\emph{IRARAH is created}}{Chapter H1 --IRARAH is created}}\label{chapter-h1-irarah-is-created}

\textbf{Scene 1}

Luca sat at a table in his small apartment. The windows were open, and the muffled sounds of the street drifted in. Printouts, notes, and old texts from various contexts lay before him. He read, highlighted, and set aside. Some pages he crumpled and threw in the trash. Others he rearranged. There was no set order that remained. He worked slowly, without a finish line.

\textbf{Scene 2}

He opened a repository on his screen. The code wasn\textquotesingle t his. He read comments, tracked changes. He added his own note, precise, concise. No explanation. He closed the window without waiting for a response.

\textbf{Scene 3}

Three people sat at a small table in the café on the corner. Cups clinked, voices overlapped. Luca said little. He asked a question, listened to the answers. No agreement was reached. Someone laughed, another changed the subject. Luca paid, stood up, and left.

\textbf{Scene 4}

Elsewhere, in a project unrelated to him, an objection arose. A short paragraph, factually worded. It was read, commented on, and moved aside. No one asked who had written it. The work continued, slightly modified.

\textbf{Scene 5}

Later, alone again, Luca typed a name into a file. No title page, no date. The name was fitting, without being intrusive. He saved it and closed the computer. Outside, it grew dark.

\subsection{\texorpdfstring{Chapter H2 --\emph{Before the workshop}}{Chapter H2 --Before the workshop}}\label{chapter-h2-before-the-workshop}

Michael walked through the university courtyard early in the morning. The paving stones were still damp from the night\textquotesingle s cleaning. He carried a folder under his arm, his steps steady. Emails, a schedule, and the workshop materials awaited him in the office. Colleagues nodded to him, briefly and respectfully. He sat down, opened his computer, and reviewed the points one last time. Everything was ready.

In Pompeii, Martina stood at the edge of an excavation site. The sun was already high, dust hung in the air. She took measurements, relayed instructions, and checked finds. Students listened, working on their own projects. During a break, she drank water and gazed at the walls that had stood for centuries. She didn\textquotesingle t give it another thought.

In Rome, Luca sat in a sparsely populated library. He moved between documents, messages, and open windows. Names appeared, then disappeared again. He wrote short messages, received replies, and saved them. There was no fixed location from which he worked, only connections.

Michael closed his office door that evening. The city was quiet. He went home, read a few more pages, and put the book down.

Martina returned to her apartment, washed the dust off her hands, and prepared something simple to eat.

Luca left the library and mingled with the people on the street.

Everything had its place. Nobody knew that it was all connected.

\section{persons}\label{persons}

\subsection{Michael Phillips}\label{michael-phillips}

Born in 1973 in Boston, Massachusetts.\\
I grew up in a teacher\textquotesingle s household.\\
Studied theology at Boston College, and physics at Boston University.\\
Entry into the Jesuit order, ordination as a priest in Boston.\\
Further studies and teaching at the Pontificia Università Gregoriana in Rome.\\
Father of Martina and Luca.

\subsection{Julia Rossi}\label{julia-rossi}

Born in Italy.\\
Studied psychology at the Pontificia Università Gregoriana in Rome.\\
Lived in Rome with Michael Phillips during his studies.\\
Martina\textquotesingle s mother.\\
After the separation, she worked as a psychologist in Pompeii.\\
Martina raised her alone.

\subsection{Martina Rossi}\label{martina-rossi}

Born in Rom.\\
Childhood and youth in Pompeii.\\
Studied archaeology in Rome.\\
Works in an archaeological context.\\
She has known her father since early years, but grows up without him.

\subsubsection{\texorpdfstring{\hfill\break
}{ }}\label{section}

\subsection{Maria}\label{maria}

Born in Italy.\\
Worked as a secretary at the Pontifical College of German and Hungarian Studies in Rome.\\
Luca\textquotesingle s mother.\\
He kept its origin secret.

\subsection{Luca}\label{luca}

Born in Rom.\\
He grew up with his mother.\\
Received a fragmentary, predominantly self-taught education.\\
Became active in digital and informal networks at an early age.\\
He maintained loose contact with peripheral church and academic communities.

\subsection{Father Matteo Conti}\label{father-matteo-conti}

Born in 1948 in Italy.\\
Jesuit, Moralphilosoph.\\
He taught at the Pontificia Università Gregoriana in Rome.\\
He was relieved of his duties there in 2005.\\
Remained present in the church community.\\
I met Luca once by chance.

\end{document}
